\documentclass[11pt]{article}
\usepackage[utf8]{inputenc}
\usepackage{amsmath,amssymb}
\usepackage{graphicx}
\usepackage{booktabs}
\usepackage{authblk}
\usepackage[margin=1in]{geometry}
\usepackage[hidelinks]{hyperref}
\usepackage{xurl}
\usepackage{setspace}

\title{RealQM vs Standard Model\\
\large Extension or Point, One Term or Many}
\author[1]{Claes Johnson}
\affil[1]{KTH Royal Institute of Technology, SE-100 44 Stockholm, Sweden \\
\texttt{claesjohnson@gmail.com}}
\date{\today}

\begin{document}
\maketitle

\begin{abstract}
RealQM and the Standard Model (SM) differ, at bottom, in a single assumption: \emph{what a physical
object fundamentally is}. For the SM the fundamental carriers are point excitations of fields --- electrons
and quarks treated as objects of zero spatial extension. For RealQM they are \emph{extended charge
densities in real three-dimensional space}, on the principle that a physical object must occupy space. We
trace this one divergence through its consequences: that extended objects on non-overlapping domains
acquire a \emph{unique identity by spatial occupation} --- absent from both standard quantum mechanics and
the SM, where identical particles are indistinguishable; what each theory takes as given; what
deep-inelastic scattering (DIS) does and does not establish; why the ``resolution'' that probes
$10^{-18}$~m is a wave and Fourier notion rather than a microscope; and where each theory actually
delivers results. The honest
conclusion is a \emph{scope split}, not a knockout. The SM rests on assumption-laden foundations --- the
point particle with its infinite self-energy, tamed only by renormalization --- yet carries an
overwhelming, interlocking record of confirmed prediction at high energy. RealQM rests on a cleaner
ontology --- extension, finite self-energy, real-space densities --- and delivers atomic and nuclear
structure from charge and size alone, but has produced nothing below the nucleon. The dividing question is
philosophical, extension versus the point; and it is worth noting that physics at its own frontier ---
fields, and strings --- already leans toward extension.
\end{abstract}

\setstretch{1.06}

\section{One assumption, not many}
\label{sec:one}

Comparisons of RealQM~\cite{RealQM} with the Standard Model usually list many differences: real space versus
configuration space, determinism versus the Born rule, non-overlap versus antisymmetry, Coulomb versus the
gauge forces. These are real, but they are downstream. Beneath them is a single assumption about which the
two frameworks disagree, and from which much of the rest follows: \emph{what kind of thing a fundamental
object is}.

The Standard Model's fundamental carriers --- electrons, quarks, and the rest --- are, in its formalism,
\emph{point} objects: zero-extension excitations of quantum fields. RealQM's carriers are \emph{extended}
objects: charge densities $\psi_i^2$ spread over regions of real three-dimensional space. Everything in what
follows is the working-out of that one difference.

\section{Extension versus the point}
\label{sec:ext}

Take the RealQM assumption first, stated as a principle: \emph{a physical object must have spatial
extension.} This is old --- Descartes' \emph{res extensa}, extension as the mark of the material --- and it
has a concrete physical edge. A charge of zero extension has \emph{infinite} electrostatic self-energy,
$\int E^2\,dV$ diverging at the origin; an extended charge density has finite self-energy. So the point is
not merely inelegant but ill-defined, and the extended density is the well-posed object. A thing that
occupies no space is a location, not an object.

The Standard Model does not deny this so much as work around it. In quantum field theory the divergence is
softened (logarithmic rather than linear for the electron) and absorbed by \emph{renormalization} --- a
formal subtraction of infinities that yields correct numbers but that Dirac, among others, regarded as a
symptom rather than a cure. So the point particle survives in the SM not because it is physically
unproblematic but because renormalization made its problems \emph{manageable}, and because no experiment
has ever \emph{required} its abandonment: no internal structure has been resolved.

Two clarifications keep this honest. First, the assumption ``objects must have extension'' is a
\emph{metaphysical choice}, adopted and then imposed on one's theories --- not a theorem physics can
derive. Second, and pointing the other way, the SM's own \emph{deepest} objects already satisfy it: the
fundamental entities of quantum field theory are \emph{fields}, extended by definition over all space, and
``particles'' are their localized excitations. The literal point is a manner of speaking about a localized
field mode. So extension is not a heterodox demand; it is arguably what physics already believes,
obscured by particle language, and made most literal --- extension in real 3D space, for any number of
particles --- by RealQM.

\section{Extension individuates: identity by spatial occupation}
\label{sec:identity}

A consequence of the extension assumption deserves separate statement, because it is a basic element of
RealQM present in \emph{neither} standard quantum mechanics nor the Standard Model. Extended objects on
non-overlapping domains are \emph{individuated by the space they occupy}: an electron, in RealQM, is
``the one filling this region,'' and two electrons are distinct because they occupy distinct territories.
Identity is conferred by spatial occupation.

This is the classical and Leibnizian notion of individuation --- impenetrability (two bodies cannot
occupy the same place) and the identity of indiscernibles (things are distinguished by their properties,
their location among them) --- restored at the quantum level. And it makes the Pauli principle not a
separate postulate but the \emph{same} fact told twice: two extended charge clouds cannot overlap, and
that single geometric fact both \emph{excludes} them and \emph{individuates} them. Exclusion and identity
are one and the same.

Standard quantum mechanics abandons this. Its identical particles are \emph{indistinguishable}: the
$N$-electron wave function lives on configuration space and is antisymmetrised under permutation, so no
electron is ``this one'' rather than ``that one.'' Pauli exclusion follows from the antisymmetry, not from
any territory; the particles carry no individual identity, and the identity of indiscernibles is,
notoriously, violated --- many electrons, all exactly alike, none of them individuated.

The Standard Model inherits the same identity-lessness, and arguably deepens it: particles are
\emph{quanta} of fields, and identical quanta are interchangeable by construction --- indistinguishability
is the very basis of the quantum statistics (Fermi, Bose) on which the theory is built. One cannot say
which quantum is which. So identity by spatial occupation is absent in the SM as well.

The contrast is therefore three-way and sharp. \emph{RealQM:} each electron has a unique identity,
conferred by the region it occupies, with exclusion and individuation the same geometric fact.
\emph{Standard QM:} no individual identity, exclusion by antisymmetry on configuration space. \emph{SM:}
no individual identity, particles as interchangeable field quanta. Only an ontology of \emph{extended}
objects occupying space can individuate by occupation; a point, which occupies no space, and a field
quantum, which is interchangeable by definition, cannot. Identity by spatial occupation is thus a direct
dividend of the extension assumption --- and one of the cleanest markers separating RealQM from both
standard formulations.

\section{What each theory takes as given}
\label{sec:given}

Every theory takes something as given and explains what is built from it. The two frameworks are strikingly
asymmetric in how much.

RealQM, in its nuclear form~\cite{RealNucleus}, takes as \emph{dimensional} input only two quantities, both
plain experimental facts that set the units: the unit \emph{charge} $\pm1$, and one \emph{length scale},
the size of the proton. Given the electric force and the variational principle, everything dimensionless
follows as prediction --- binding-energy ratios, the shape of the binding-per-nucleon curve, saturation.
The proton's femtometre size is \emph{measured}, not derived; it enters exactly as the Bohr radius enters
atomic physics, as a unit-setting scale. Charge and extension are about the most primitive properties an
object can have.

The Standard Model takes as given a far larger apparatus: a gauge group, a set of fields, roughly nineteen
free parameters (masses, mixings, couplings), and the machinery of renormalization. In return it explains
an enormous range of phenomena RealQM does not touch. The asymmetry is real and cuts both ways: RealQM
assumes very little and explains a bounded domain; the SM assumes a great deal and explains a vast one.

\section{The Lagrangian compared: three forces, or one operator}
\label{sec:lagrangian}

The asymmetry shows most starkly in the equations themselves. The Standard Model is a relativistic gauge
quantum field theory with gauge group $SU(3)_c\times SU(2)_L\times U(1)_Y$ and Lagrangian
\begin{equation}
\begin{aligned}
\mathcal L_{\rm SM}=\;&\underbrace{-\tfrac14 G^A_{\mu\nu}G^{A\mu\nu}}_{\text{strong}}
\;\underbrace{-\tfrac14 W^a_{\mu\nu}W^{a\mu\nu}-\tfrac14 B_{\mu\nu}B^{\mu\nu}}_{\text{electroweak}}
\;+\;\underbrace{\sum_\psi\bar\psi\,i\gamma^\mu D_\mu\psi}_{\text{fermions + couplings}}\\[3pt]
&+\;\underbrace{(D_\mu\Phi)^\dagger(D^\mu\Phi)-V(\Phi)}_{\text{Higgs}}
\;-\;\underbrace{\big(\textstyle\sum y\,\bar\psi_L\Phi\,\psi_R+\text{h.c.}\big)}_{\text{Yukawa}},
\end{aligned}
\label{eq:SM}
\end{equation}
the covariant derivative $D_\mu=\partial_\mu-ig_s G-ig\,W-ig'B$ coupling every fermion to all three forces
at once, with of order nineteen free parameters. Keeping \emph{only} the electromagnetic term left
unbroken after $SU(2)_L\times U(1)_Y\to U(1)_{\rm em}$ gives quantum electrodynamics,
\begin{equation}
\mathcal L_{\rm QED}=-\tfrac14 F_{\mu\nu}F^{\mu\nu}+\bar\psi\big(i\gamma^\mu\partial_\mu-m\big)\psi
+e\,q\,\bar\psi\gamma^\mu\psi\,A_\mu .
\label{eq:QED}
\end{equation}
In the non-relativistic, electrostatic limit --- Dirac $\to$ Schr\"odinger, and the photon reduced to its
scalar potential $A_0=\varphi$ with $-\nabla^2\varphi=4\pi\rho$, i.e.\ instantaneous Coulomb --- this
becomes the multiphase Schr\"odinger Lagrangian
\begin{equation}
\mathcal L_{\rm RealQM}=\sum_i\int_{\Omega_i}\!\Big[\tfrac{i}{2}\big(\psi_i^*\dot\psi_i-\dot\psi_i^*\psi_i\big)
-\tfrac{1}{2m_i}|\nabla\psi_i|^2\Big]-\tfrac12\!\iint\frac{\rho(x)\rho(y)}{|x-y|},\qquad
\rho=\sum_i q_i|\psi_i|^2,
\label{eq:LRealQM}
\end{equation}
read now in \emph{real three-dimensional space} on non-overlapping domains rather than on configuration
space. Being first order in time, its Hamiltonian --- the conserved energy the ground state minimises ---
is the energy functional
\begin{equation}
E=\sum_i\tfrac12\!\int_{\Omega_i}|\nabla\psi_i|^2\;+\;\tfrac12\!\iint\frac{\rho(x)\rho(y)}{|x-y|},\qquad
\rho=\sum_i q_i\psi_i^2,\quad q_i=\pm1.
\label{eq:Efunc}
\end{equation}
So the Lagrangian \eqref{eq:LRealQM} yields the energy functional \eqref{eq:Efunc} as its Hamiltonian, and
RealQM is the variational problem for the latter.

In this precise sense RealQM is \emph{one term of the Standard Model} --- its electromagnetic term ---
kept alone and read as extended charge densities. The electromagnetism is included in full, in two stages:
the Coulomb part \eqref{eq:Efunc} carries the ground-state binding, while the transverse photon --- and
hence \emph{currents, magnetism, and radiation} --- is restored by the time-dependent form
\eqref{eq:LRealQM}, whose current is $\mathbf J=\frac{q}{m}\mathrm{Im}(\psi^*\nabla\psi)$. The strong
force, the weak force, the Higgs and the gauge apparatus are simply absent, and RealNucleus is the claim
that this one electromagnetic term suffices for what the SM assigns to three: binding (strong) and
instability (weak) both follow from Coulomb alone.

\begin{center}
\small
\renewcommand{\arraystretch}{1.2}
\begin{tabular}{lll}
\hline
 & \textbf{Standard Model} & \textbf{RealQM} \\
\hline
strong / weak / Higgs / Yukawa & present & dropped \\
electromagnetic term & present & kept (Coulomb $+$ full EM in the time-dependent form) \\
matter & Dirac fields, point quanta & charge densities, extended \\
space & $3N$ configuration space & real $\mathbb R^3$ \\
exclusion & antisymmetry & non-overlap of domains \\
free parameters & ${\sim}19$ & charge and one length scale \\
\hline
\end{tabular}
\end{center}

\paragraph{One operator: the Laplacian and its inverse.} A deeper unity underlies \eqref{eq:Efunc}, and it
bears directly on extension. Both terms are built from a \emph{single} operator --- the Laplacian
$\nabla^2$ of three-dimensional space. The kinetic term is the Laplacian acting directly: varying
$\tfrac12\int|\nabla\psi|^2$ gives $-\nabla^2\psi$. The Coulomb term is the Laplacian \emph{inverted}: the
potential $\varphi(x)=\int\rho(y)/|x-y|\,dy$ solves Poisson's equation $-\nabla^2\varphi=4\pi\rho$, so the
kernel $1/|x-y|$ is the Green's function of $-\nabla^2$. Kinetic energy is thus
$\langle\psi,-\nabla^2\psi\rangle$ and Coulomb energy is $\langle\rho,(-\nabla^2)^{-1}\rho\rangle$: the
operator and its inverse, one local (differential), one nonlocal (integral), both the same Laplacian of
real space. The whole theory is the Laplacian, read forwards and backwards.

\paragraph{Kinetic energy creates extension.} This clarifies what the kinetic term \emph{is}. It is not
motion --- nothing in a RealQM ground state moves. It is the internal pressure that \emph{creates spatial
extension}: Coulomb attraction alone would collapse a charge cloud onto the kernel, and the gradient term
is what resists the collapse, fixing a finite size (their balance sets the Bohr radius, and at nuclear
scale the femtometre). Read this way the kinetic term is a form of \emph{controlled self-repulsion} --- the
``self-Coulomb'' of each charge cloud on its own domain, the density's resistance to being squeezed toward
a point. It is repulsive, like the Coulomb term, and of the same Laplacian form, differing only in that the
\emph{mutual} interaction is the nonlocal $1/r$ while the \emph{self}-repulsion is its local differential
counterpart. (This is also why the inter-electron Coulomb sum correctly omits the diagonal $j=i$: the
self-interaction is already supplied, in local form, by the gradient term.) So extension is not merely
assumed in RealQM but \emph{generated} --- by a kinetic term that is itself a controlled electrostatic
self-repulsion --- and the entire dynamics reduces to one operator, the Laplacian, and its inverse. Against
the Standard Model's three forces and nineteen parameters, this is as parsimonious as a physical theory can
be.

\section{What deep-inelastic scattering establishes --- and what it interprets}
\label{sec:dis}

The sharpest test of the two ontologies is deep-inelastic scattering, standardly read as showing that the
proton is made of point-like quarks. It is essential to separate what DIS \emph{measures} from what it
\emph{interprets}.

\emph{Measured, and not in dispute:} at high momentum transfer the proton scatters as though built of
constituents that show no resolved structure and no scale --- Bjorken scaling. This is data. The proton is
also, incontestably, a definite object: measured mass, charge, radius ($\approx 0.84$~fm), and magnetic
moment, on effectively free protons (the atomic binding, ${\sim}$eV, is negligible against the GeV probe).
The target exists; the structure functions are real.

\emph{Interpreted, and model-laden:} that the constituents are \emph{literal zero-extension points} is not
data but the parton-model reading, resting on the impulse approximation. By the extension principle of
\S\ref{sec:ext} such a point is not a physical object; the honest translation of ``pointlike'' is
\emph{structureless down to current resolution}, i.e.\ a very concentrated \emph{extended} object, not a
mathematical point. And QFT itself does not require the bare point: the quark is an excitation of an
extended field.

RealQM neither disputes the DIS data nor, at present, reproduces it. Its position is scope: it is a theory
of nuclear \emph{binding}, taking the proton as a charge unit of given size, and it says nothing about the
proton's internal structure --- a separate, higher-energy question it neither needs nor contests. The
extended-object reading of DIS (concentrated constituents, unresolved below $10^{-18}$~m) is entirely
compatible with the data; a RealQM-style \emph{computation} of the proton structure functions, however, has
not been done, and the confinement of such constituents is not something Coulomb alone supplies.

\section{Resolution is a wave notion, not a microscope}
\label{sec:res}

That experiments ``resolve $10^{-18}$~m'' can seem to conflate energy with spatial resolution. It does not
--- but the link is not self-evident, and it rests on the wave nature of matter, which is precisely the
ontology RealQM shares. Two steps: de~Broglie ($\lambda=h/p$, a fast probe is a short wave) and the fact
that waves resolve structure only at their wavelength (a long wave washes over small features; a short one
feels them). More precisely, position and momentum are Fourier conjugates ($\Delta x\,\Delta p\sim\hbar$),
and scattering measures the \emph{form factor} --- the Fourier transform of the charge distribution --- so
sampling short-distance structure requires large momentum transfer, $d\sim\hbar/q$. This is not imaging: no
picture is formed, and ``resolution'' is shorthand for the length scale to which the momentum transfer is
sensitive. Crucially, the whole relationship is a \emph{wave} statement --- momentum as spatial frequency,
structure read through Fourier transforms --- and thus sits inside RealQM's own picture rather than against
it. High energy accesses fine structure \emph{because} matter is waves; take away the wave, and the
inference collapses.

\section{The honest scorecard: scope, not knockout}
\label{sec:score}

It is tempting, having pressed the SM's foundations, to declare its case weak. That would confuse two
different games.

\emph{On foundations}, the SM is genuinely vulnerable: the literal point is indefensible, the
renormalization that saves it is a formal subtraction Dirac called a symptom, the celebrated precision
tests lean on measured inputs and are limited by uncertainties (the ``twelve digits'' overstate the
independent content), and the interpretation of DIS is model-laden. On this front RealQM's principle ---
extension --- is the sounder one.

\emph{On the predictive record}, the SM's cards are not weak; they are the strongest in science: the hadron
spectrum, jets, the running coupling, the $W$, $Z$, and Higgs, cross-sections and lifetimes confirmed across
many orders of magnitude, an interlocking web that no rival matches. Exposing shaky \emph{foundations} is
not the same as showing the \emph{predictions} fail; it does not.

And fairness cuts the same knife at RealQM. Below the nucleon its cards are not weak but \emph{empty}: no
quark spectrum, no structure functions, no confinement mechanism, no results at all. Its strength is at its
\emph{own} scale --- atoms, molecules, chemistry, and nuclear binding from charge and size --- where it
delivers what the SM's high-energy apparatus is not built to compute cheaply, and where, in the nuclear
case, QCD has after five decades produced no parameter-free binding energy at all~\cite{QCDBinding}.

The honest summary is therefore a table of \emph{scope}, not a verdict of victory:

\begin{center}
\begin{tabular}{@{}p{0.31\linewidth} p{0.31\linewidth} p{0.31\linewidth}@{}}
\toprule
 & \textbf{Standard Model} & \textbf{RealQM} \\
\midrule
fundamental object & point excitation of a field & extended charge density in real 3D \\[3pt]
self-energy of a charge & infinite, renormalised & finite \\[3pt]
inputs & gauge group, ${\sim}19$ parameters, renormalisation & charge and one length scale \\[3pt]
foundations & assumption-laden (the point) & cleaner (extension) \\[3pt]
predictive record & overwhelming, high-energy & bounded, atomic/nuclear \\[3pt]
sub-nucleon regime & its home ground & no results \\[3pt]
atomic/nuclear binding & costly; QCD delivers no parameter-free nuclear binding & delivered from one scale \\
\bottomrule
\end{tabular}
\end{center}

\section{Conclusion: the dividing assumption, and where physics already leans}
\label{sec:concl}

The disagreement reduces to one assumption: is the fundamental object a point or an extended thing? Stated
plainly, the point loses on physical grounds --- infinite self-energy, an object with no extension is not an
object --- and RealQM is the framework that takes the alternative most literally, as extended charge
densities in real three-dimensional space, with finite self-energy and, at the atomic and nuclear scale,
concrete computable results. Stated just as plainly, that philosophical win does not touch the Standard
Model's predictive record, which remains the strongest in science, nor does it fill RealQM's empty hand
below the nucleon.

What it does establish is where each theory rightly stands. The Standard Model owns the high-energy,
sub-nucleon world by a vast confirmed record built on a foundation it does not examine. RealQM owns the
atomic and nuclear world by a foundation --- extension, charge, size --- it examines closely and computes
from directly. And the choice between a point and an extended object, far from being RealQM's private
heterodoxy, is one physics has already made at its own frontier, where the fundamental objects are
extended fields and extended strings. RealQM's distinctive claim is only to make that same choice
\emph{concrete and calculable} in real three-dimensional space --- and to show that, at its scale, charge
and extension are enough.

\begin{thebibliography}{9}
\bibitem{RealQM} C.~Johnson, \emph{Real Quantum Mechanics as Multiphase 3D Continuum Mechanics},
manuscript, 2026; \url{https://claes542.github.io/RealMolecule/gallery.html}.
\bibitem{RealNucleus} C.~Johnson, \emph{RealNucleus: Nuclear Binding as Dual Coulomb Confinement, without
a Strong or Weak Force}, 2026.
\bibitem{QCDBinding} C.~Johnson, \emph{The Binding Energy Problem QCD Never Solved --- and a Coulomb Model
That Does}, 2026.
\end{thebibliography}

\end{document}
